%! Unit 1: Functions and Change — Summative Assessment — Practice Test --- Study Copy (KEY)
\documentclass[10pt]{article}
\usepackage{precalculus-article}
\usepackage{precalculus-key}   % pulls in -boxes; do NOT also load -boxes

% Part-divider strip
\newcommand{\parthead}[1]{%
  \vspace{6pt}\noindent
  \begin{headlinebox}{plum}{\color{white}\bfseries #1}\end{headlinebox}%
  \vspace{4pt}}

\begin{document}

\pageheader{Unit 1: Functions and Change}{Practice Test --- Study Copy --- Answer Key}
\namedateperiod

\vspace{0.10in}
\begin{remindbox}
\small
\textbf{This is a practice test.} It mirrors the real Unit 1 test in length, format,
and the ideas it covers, but the numbers and contexts differ. Work every problem
with no notes, then check your answers against the key.
\end{remindbox}

% ======================================================================
\parthead{Part A --- Vocabulary (8 pts)}
Write the letter of the best description next to each term.

\vspace{4pt}
\noindent
\begin{minipage}[t]{0.42\textwidth}
\begin{enumerate}[leftmargin=2.2em, itemsep=7pt, label=\arabic*.]
  \item Function \ans{C}
  \item Domain \ans{G}
  \item Range \ans{A}
  \item Average rate of change \ans{F}
  \item Concave up \ans{H}
  \item Parent function \ans{B}
  \item Composition \ans{E}
  \item Inverse function \ans{D}
\end{enumerate}
\end{minipage}%
\hfill
\begin{minipage}[t]{0.55\textwidth}
\small
\begin{description}[leftmargin=1.4em, itemsep=3pt]
  \item[(A)] The set of every output the function actually produces.
  \item[(B)] The simplest form of a family, before any shift or stretch.
  \item[(C)] A rule that sends each input to \emph{exactly one} output.
  \item[(D)] The function that reverses the input--output pairs of $f$.
  \item[(E)] Feeding the output of one function into another: $f(g(x))$.
  \item[(F)] The slope of the secant line, $\dfrac{f(b)-f(a)}{b-a}$.
  \item[(G)] The set of every input the function is allowed to take.
  \item[(H)] Bending upward --- the rate of change is increasing.
\end{description}
\end{minipage}

% ======================================================================
\parthead{Part B --- Multiple Choice (12 pts)}
Circle the best answer.

\begin{enumerate}[leftmargin=1.9em, itemsep=9pt]
  \item A relation fails to be a function exactly when
    \hfill (A) two inputs share an output \quad \textbf{(B) one input has two outputs} \ans{$\leftarrow$} \\
    \phantom{x}\hfill (C) the graph crosses the $x$-axis \quad (D) the domain is restricted
  \item If a function is decreasing and concave up on an interval, its output is
    \hfill (A) falling faster and faster \quad \textbf{(B) falling, but more and more slowly} \ans{$\leftarrow$} \\
    \phantom{x}\hfill (C) rising more and more slowly \quad (D) changing at a constant rate
  \item A table of $f$ has equally spaced inputs and \textbf{constant second} differences.
    The best model is
    \hfill (A) linear \quad \textbf{(B) quadratic} \ans{$\leftarrow$} \quad (C) exponential \quad (D) none of these
  \item The graph of $g(x) = f(x+3)$ is the graph of $f$ shifted
    \hfill (A) $3$ right \quad \textbf{(B) $3$ left} \ans{$\leftarrow$} \quad (C) $3$ up \quad (D) $3$ down
  \item If $f(x) = x^2$ and $g(x) = x - 1$, then $(f \circ g)(4)$ equals
    \hfill \textbf{(A) $9$} \ans{$\leftarrow$} \quad (B) $15$ \quad (C) $16$ \quad (D) $3$
  \item A function has an inverse that is also a function exactly when the function is
    \hfill (A) increasing somewhere \quad \textbf{(B) one-to-one} \ans{$\leftarrow$} \quad
    (C) concave up \quad (D) defined at $0$
\end{enumerate}

\boxguard[8]
\begin{teachernote}
Item 1: students who pick (A) are confusing the vertical line test with the horizontal
one --- two inputs \emph{may} share an output. Item 4: the inside change $x+3$ moves the
graph the ``wrong'' way; expect (A) as the common miss. Item 5: watch for students who
compute $g(f(4)) = 15$ instead of $f(g(4)) = f(3) = 9$.
\end{teachernote}

% ======================================================================
\parthead{Part C --- Short Answer \& Computation (30 pts)}
Show your work.

\begin{enumerate}[leftmargin=1.9em, itemsep=12pt]
  \item Let $f(x) = 2x^2 - 5$. Compute \textbf{(a)} $f(3)$, \quad \textbf{(b)} $f(-2)$, \quad
        \textbf{(c)} the value of $x$ (there are two) for which $f(x) = 3$.
        \ans{(a) $f(3) = 2(9) - 5 = 13$. \ (b) $f(-2) = 2(4) - 5 = 3$. \ (c) $2x^2 - 5 = 3 \Rightarrow 2x^2 = 8 \Rightarrow x^2 = 4 \Rightarrow x = \pm 2$.}
  \item The graph of $y = f(x)$ is shown below.
        \textbf{(a)} On what interval is $f$ decreasing?
        \textbf{(b)} Is $f$ concave up or concave down on the whole picture?
        \textbf{(c)} List the zeros of $f$.
        \textbf{(d)} What is $f(3)$?

        \vspace{2pt}
        \begin{center}
        \begin{tikzpicture}[scale=0.68]
          \draw[linegray, very thin] (-0.4,-2.4) grid (6.4,3.4);
          \draw[->] (-0.4,0) -- (6.6,0) node[right]{\small $x$};
          \draw[->] (0,-2.4) -- (0,3.6) node[above]{\small $y$};
          \foreach \x in {1,...,6} \draw (\x,0.09) -- (\x,-0.09) node[below]{\scriptsize $\x$};
          \foreach \y in {-2,-1,1,2,3} \draw (0.09,\y) -- (-0.09,\y) node[left]{\scriptsize $\y$};
          \draw[plum, very thick] plot[smooth] coordinates
            {(0,2.5) (1,0) (2,-1.5) (3,-2) (4,-1.5) (5,0) (6,2.5)};
        \end{tikzpicture}
        \end{center}
        \ans{(a) Decreasing on $0 < x < 3$ (increasing on $3 < x < 6$). \ (b) Concave up everywhere shown --- the curve bends upward and the rate of change is increasing. \ (c) Zeros at $x = 1$ and $x = 5$. \ (d) $f(3) = -2$ (the low point).}
  \item A tank is draining. The volume $V$ (gallons) after $t$ minutes is given below.

        \vspace{2pt}
        \begin{center}
        \renewcommand{\arraystretch}{1.3}
        \begin{tabular}{|c|c|c|c|c|}
        \hline
        $t$ (min) & $0$ & $4$ & $8$ & $12$ \\ \hline
        $V$ (gal) & $180$ & $148$ & $124$ & $108$ \\ \hline
        \end{tabular}
        \end{center}
        \vspace{2pt}
        \textbf{(a)} Find the average rate of change of $V$ from $t = 0$ to $t = 8$.
        \textbf{(b)} Give its units and say in one sentence what it means for the tank.
        \textbf{(c)} Is the tank draining faster or slower as time goes on? Justify with the numbers.
        \ans{(a) $\dfrac{V(8) - V(0)}{8 - 0} = \dfrac{124 - 180}{8} = \dfrac{-56}{8} = -7$. \ (b) Gallons per minute: over the first $8$ minutes the tank lost water at an average of $7$ gallons each minute (negative because the volume is falling). \ (c) Slower. The change over each $4$-minute block is $-32$, then $-24$, then $-16$ gallons --- the drops are shrinking, so the draining is slowing down.}
  \item The table below has equally spaced inputs. Compute the first differences and the
        second differences, then name the function type (linear, quadratic, or neither).

        \vspace{2pt}
        \begin{center}
        \renewcommand{\arraystretch}{1.3}
        \begin{tabular}{|c|c|c|c|c|c|}
        \hline
        $x$    & $0$ & $1$ & $2$ & $3$ & $4$ \\ \hline
        $g(x)$ & $2$ & $5$ & $12$ & $23$ & $38$ \\ \hline
        \end{tabular}
        \end{center}
        \ans{First differences: $3,\ 7,\ 11,\ 15$ (not constant, so not linear). Second differences: $4,\ 4,\ 4$ --- constant. The function is \emph{quadratic}.}
  \item The parent function is $f(x) = \sqrt{x}$, with domain $[0,\infty)$ and range
        $[0,\infty)$. Let $g(x) = -2f(x - 1) + 4$.
        \textbf{(a)} Describe every transformation, in order.
        \textbf{(b)} State the domain and range of $g$.
        \ans{(a) $g(x) = -2\sqrt{x-1} + 4$: shift right $1$, then vertical stretch by a factor of $2$, then reflect over the $x$-axis, then shift up $4$. \ (b) Domain: $x - 1 \ge 0$, so $[1, \infty)$. Range: $\sqrt{x-1} \ge 0$, so $-2\sqrt{x-1} \le 0$ and $g(x) \le 4$ --- the range is $(-\infty, 4]$.}
  \item Let $f(x) = 3x + 1$ and $g(x) = x^2 - 2$.
        \textbf{(a)} Compute $f(g(2))$. \quad \textbf{(b)} Compute $g(f(2))$.
        \textbf{(c)} Write a formula for $(f \circ g)(x)$.
        \textbf{(d)} What do (a) and (b) show about the order of composition?
        \ans{(a) $g(2) = 4 - 2 = 2$, so $f(g(2)) = f(2) = 7$. \ (b) $f(2) = 7$, so $g(f(2)) = 49 - 2 = 47$. \ (c) $(f \circ g)(x) = 3(x^2 - 2) + 1 = 3x^2 - 5$. \ (d) $7 \ne 47$: order matters. Composition is not commutative --- $f(g(x))$ and $g(f(x))$ are generally different functions.}
  \item Let $f(x) = \dfrac{x - 5}{2}$.
        \textbf{(a)} Find $f^{-1}(x)$ using the swap-and-solve method.
        \textbf{(b)} Verify your answer by computing $f\!\left(f^{-1}(x)\right)$.
        \ans{(a) Write $y = \dfrac{x-5}{2}$, swap: $x = \dfrac{y-5}{2}$, solve: $2x = y - 5 \Rightarrow y = 2x + 5$. So $f^{-1}(x) = 2x + 5$. \ (b) $f\!\left(f^{-1}(x)\right) = \dfrac{(2x+5) - 5}{2} = \dfrac{2x}{2} = x$. $\checkmark$}
\end{enumerate}

% ======================================================================
\parthead{Part D --- Extended Response (10 pts)}
Explain your reasoning in full sentences.

\begin{enumerate}[leftmargin=1.9em, itemsep=12pt]
  \item A rectangular dog run is built against a barn wall, so only \textbf{three} sides need
        fencing. The builder has $60$ feet of fence. Let $x$ be the length of each of the two
        sides perpendicular to the barn.
        \textbf{(a)} Write the area $A$ as a function of $x$.
        \textbf{(b)} What type of function is $A$, and how do you know?
        \textbf{(c)} State the domain restriction that the context forces, and explain why.
        \textbf{(d)} State one assumption this model rests on.
        \ans{(a) The two perpendicular sides use $2x$ feet, so the side parallel to the barn is $60 - 2x$. Then $A(x) = x(60 - 2x) = 60x - 2x^2$. \ (b) Quadratic --- the highest power of $x$ is $2$, and equally spaced $x$-values would give constant second differences. It opens downward, so it has a maximum. \ (c) $0 < x < 30$: a side length must be positive, and the barn side $60 - 2x$ must also be positive, which forces $x < 30$. Outside that window the ``rectangle'' has no meaning. \ (d) Any one of: no fence is used along the barn; the barn wall is at least $60 - 2x$ feet long; the run is exactly rectangular; no fence is lost to a gate or to overlap.}
  \item Let $f(x) = x^2$ with domain all real numbers.
        \textbf{(a)} Explain why $f$ has no inverse function on that domain.
        \textbf{(b)} Give a domain restriction that fixes the problem, and write the resulting
        inverse.
        \textbf{(c)} Explain how the graph of a function and the graph of its inverse are related.
        \ans{(a) $f$ is not one-to-one: $f(2) = f(-2) = 4$. Reversing the pairs would send the input $4$ to two different outputs, $2$ and $-2$, and that is not a function. (Equivalently, $f$ fails the horizontal line test.) \ (b) Restrict the domain to $x \ge 0$. On $[0,\infty)$ the function is one-to-one, and $f^{-1}(x) = \sqrt{x}$ with domain $[0,\infty)$. \ (c) They are reflections of each other over the line $y = x$: the point $(a,b)$ on $f$ becomes $(b,a)$ on $f^{-1}$, because the inverse swaps input and output.}
\end{enumerate}

\boxguard[8]
\begin{teachernote}
\textbf{Scoring Part D (5 pts each).} Award full credit only when the student both gives
the correct result and \emph{justifies} it. D1: 2 pts (a) $A(x) = 60x - 2x^2$ with the
$60 - 2x$ side reasoned out; 1 pt (b) quadratic \emph{with} a reason; 1 pt (c) $0 < x < 30$
tied to the context, not just stated; 1 pt (d) any defensible assumption. A student who
writes $2x + 2y = 60$ (all four sides) has missed the barn --- give partial credit for
correct method on a wrong perimeter. D2: 2 pts (a) not one-to-one \emph{with} a
counterexample pair; 2 pts (b) restriction \emph{and} $f^{-1}(x) = \sqrt{x}$; 1 pt (c)
reflection over $y = x$. Full-credit language for the arc: \emph{one-to-one $\to$ reverse
the pairs $\to$ restrict to make it work $\to$ reflect over $y = x$}.
\end{teachernote}

\end{document}
